% SPDX-FileCopyrightText: 2026 Hasan Melih Akbulut
% SPDX-License-Identifier: CC-BY-4.0
%
% Chapter 2: the architecture of the system-on-chip. Every number in
% this file is cited to a document or a source file in this repository.

\chapter{Architecture}
\label{ch:arch}

The part is \code{soc\_top}: a 32-bit RISC-V management core, a system
fabric, a memory subsystem with a code over every stored bit, a
core-local interruptor, a general-purpose timer carrying a three-stage
watchdog, a boot controller, two telemetry blocks, a console UART, a
GPIO port, a QSPI flash controller, two plug-and-play device tables and a
spiking-neural-network node built around a frozen shuttle die
(\dcite{60}{}, revision 0.2, 14 September 2026). Its datasheet gives its
status in three words: preliminary, pre-silicon, pre-sign-off.

One sentence travels with the rest of this chapter. \textbf{This design
publishes no clock frequency}, by a rule adopted in
\dcite{53}{section 9.2} and restated in \dcite{60}{section 0}: the
20\,ns period below is the constraint the synthesis and place-and-route
flows were run against, and every slack here is measured against it. It
is not converted into a rating and no rating is implied.

There is no block diagram in this chapter: \cref{fig:arch-panels} shows
where the placer actually put each block, and the floorplan and the
density map belong to the physical-design chapter.

\begin{figure}[tbp]
  \includegraphics[width=\linewidth]{figures/layout_panels.pdf}
  \caption{Where each block landed. Nine small multiples, one per block,
  each drawing that block's standard cells against a grey die with the
  eight vendor memory macros outlined. Hierarchy barely survives this
  flow's synthesis: of the 168,592 instances in the sign-off run's DEF
  --- 168,584 standard cells and the eight macros --- only 1{,}140 still
  spell a dot in their name, 632 of them the RAM's, and by the script's
  own prefix test exactly one carries a top-level block name. The rest
  are attributed by connectivity --- nets keep
  their hierarchical names, each anonymous cell takes the block most of
  its nets vote for, and two rounds of neighbour propagation follow
  (\code{thesis/figures/make\_layout\_figures.py}). The method leaves a
  remainder: \textbf{5{,}916 of the 60{,}694 logic cells, 9.75\,\%, are
  unattributed} and are drawn grey rather than hidden
  (\code{thesis/figures/layout\_attribution.json}). Nothing here was placement-constrained except the macros, so the
clusters are the placer's.}
  \label{fig:arch-panels}
\end{figure}

\section{The management core}

\begin{sloppypar}
The core is lowRISC Ibex in the \code{small-pmp} configuration: RV32IMC
with physical memory protection at four regions, a two-stage pipeline, a
flip-flop register file, no instruction cache, no branch-target ALU, no
writeback stage, no branch predictor. \code{hw/soc/rtl/soc\_top.v} sets
\code{BaseIsa} = 0 (RV32I), \code{RV32M} = 2 (RV32MFast),
\code{RV32B} = 0, \code{RV32ZC} = 0 (Zca), \code{RegFile} = 0
(flip-flop), \code{PMPEnable} = 1, \code{PMPNumRegions} = 4,
\code{PMPGranularity} = 0, \code{MHPMCounterNum} = 0 and
\code{SecureIbex} = 0.
\end{sloppypar}

\begin{sloppypar}
The core is vendored pristine at commit
\code{34b0705760ef3dfa00e99637432473d2be8f22f3}, Apache-2.0, and
converted to Verilog-2005 by sv2v: 30 of 30 RTL modules and 8 of 8
vendored primitives, 15{,}873 lines, zero errors and \textbf{zero patches
to Ibex} (\dcite{38}{sections 2, 4.2}). \textbf{No equivalence between
the SystemVerilog and the converted Verilog was proven}, because
upstream's own equivalence script drives a proprietary tool this project
does not have. Two corroborations were run instead --- two independent SystemVerilog
front-ends landing within \textbf{0.89\,\% of cell area and two
flip-flops}, and a self-checking program executing in Icarus --- and
\dcite{38}{section 5} refuses to call either a proof.
\end{sloppypar}

\begin{sloppypar}
PMP was not taken for free. \dcite{09}{part B track 3} had already chosen
a software architecture whose isolation mechanism is PMP, so a core
without it does not implement the design already decided. It costs
\textbf{43{,}634\,$\mu$m$^2$, $+18.8$\,\%} and \textbf{1.21\,ns} of
pre-layout closing period, because the address-match logic is replicated
per region and sits in the load/store path
(\dcite{38}{sections 8.4, 8.6}).
\end{sloppypar}

\subsection{Why there is no lockstep}

\code{SecureIbex} is 0: no shadow core, no output comparison, no dummy
instructions, no reset-all flops, no memory-interface ECC. In this Ibex
those are not separable. \code{ibex\_top.sv} lines 212--216 make
\code{Lockstep}, \code{ResetAll}, \code{DummyInstructions} and
\code{RegFileLockstepECC} localparams of \code{SecureIbex}, and
\code{MemECC} defaults to it at line 41 (\dcite{38}{section 3.3}). It is
one switch, and it was priced before it was thrown --- each
configuration synthesised at its own closing period
(\dcite{38}{section 8.4}):

\begin{center}\small
\begin{tabular}{lrrrr}
\toprule
Configuration & Cells & Flip-flops & Area, $\mu$m$^2$ & kGE \\
\midrule
\code{small}, no PMP        & 15{,}736 & 1{,}950 & 232{,}049 & 31.97 \\
\code{small-pmp}, chosen    & 18{,}921 & 2{,}106 & 275{,}683 & 37.99 \\
\code{small-pmp-sec}        & 40{,}236 & 4{,}975 & 585{,}233 & 80.64 \\
\bottomrule
\end{tabular}
\end{center}

Lockstep is \textbf{$+309{,}550\,\mu$m$^2$, $+112$\,\%} over the chosen
core and costs almost nothing in closing period, 0.27\,ns, because the
shadow core is parallel rather than in series (\dcite{38}{sections 8.5,
8.6}).

Area is not what decided it. Lockstep \emph{detects} where this project's voters \emph{mask} (this
Ibex's lockstep is one shadow core and a comparison; a triple core
lock-step votes and therefore masks as well~\cite{tcls2019}); the two
cover disjoint structures, so their costs add rather than share; and the one real overlap, ECC
over the register file, duplicates a codec this project already owns.
\code{MemECC} follows \code{SecureIbex}, so the instruction and data
buses would become 39 bits wide all the way to memory with the check bits
\emph{stored} rather than regenerated: a 22\,\% SRAM overhead on every
byte the CPU touches. \textbf{The owner chose \code{small-pmp} on
2026-08-31} (\dcite{38}{section 10 item 4}), and the same section records
what the choice takes on: the management processor becomes the least
hardened block in the design, its own fault-injection campaign stops
being deferred work, a recovery policy becomes required rather than
optional, and the decision is cheap to reverse only before floorplanning.

\begin{sloppypar}
What has happened since: the campaign was run, the register file gained a
(39,32) SECDED codec with a background scrub at
$+38{,}050.5762\,\mu$m$^2$ and $+222$ flip-flops (\dcite{43}{}), and the
watchdog is the recovery policy. The core as shipped is \textbf{22{,}338 cells, 2{,}328
flip-flops, 316{,}051.6590\,$\mu$m$^2$} (\dcite{45}{section 6.2});
everything in Ibex except the register file remains unprotected, by
decision (\dcite{60}{section 9.8}).
\end{sloppypar}

\section{The bus fabric}

Ibex presents two masters, an instruction port and a data port, each
speaking a simple request/grant/response protocol; it has no AHB, no AXI
and no burst. The interconnect takes that protocol as the internal fabric
and puts \emph{real} AMBA 3 APB at the peripheral boundary.
\textbf{There is no AHB in this SoC and none is claimed}
(\dcite{39}{sections 3, 4, 7}).

The obvious argument, area, is the one the record refuses. An honest AHB-Lite cost probe was built and synthesised at \textbf{17{,}154.32\,$\mu$m$^2$ against
14{,}807.32}, $+15.9$\,\% on the fabric and \textbf{$+0.85$\,\% of the
CPU}. Nine tenths of one per cent is not a reason to reject a bus; what
the measurement is good for is excluding the opposite mistake, that the
decision was made by picking the smaller thing (\dcite{39}{section 6.3}).

Five arguments decided it. \dcite{08}{section 3} defines proximity to
the reference family at the software-visible layer --- addresses,
register maps, slot geometry, discovery tables --- and the wires between
core and RAM are not in it. Not one peripheral core the IP survey
selected is AHB-native, so every one needs a wrapper whichever bus is
chosen, and an APB wrapper is the smaller one. \dcite{08}{row 1}'s
``single AHB-lite layer'' with three masters is a category error, because
\textbf{AHB-Lite has exactly one master by definition} and Ibex alone
supplies two. And AHB folds ``your request was accepted'' and ``your
response is here'' onto one wire, so an arbiter that back-pressures a
master simultaneously tells it its outstanding data phase has not
completed; the probe parks the response instead, and \textbf{67 of the
arbiter's 71 flip-flops are parking registers}.

Evidence settles the rest: nothing in the pinned toolchain is an AHB
protocol checker, while APB's eight-odd rules are proved by k-induction
against ghost registers independent of the bridge's own storage
(\dcite{39}{section 5.3}).

\code{soc\_bus.v} carries two masters onto \textbf{six slave ports} ---
\code{ram}, \code{rom}, \code{clint}, \code{npu}, \code{apb},
\code{pnp} --- plus an internal error slave (\dcite{60}{section 5.2}).
Three of its properties are decisions. It allows \textbf{two outstanding
requests per master}, because Ibex's prefetch buffer issues two and one
would halve instruction-fetch bandwidth. \textbf{All of a master's
outstanding requests must go to one slave}, because Ibex requires
in-order responses and slaves differ in latency; the alternative was a
reorder buffer that must store a whole response, and the restriction
costs one dead cycle when a master changes region. And arbitration is
\textbf{round-robin}, which makes no-starvation a property of this file
rather than an argument about the core: the measured worst gap between
grants is 2 cycles against an asserted bound of 8.

An unmapped access is a \textbf{bus error, never a read of zero}: the
error slave is inside the fabric, always ready, and Ibex turns its
response into an access fault. Properties F1--F9 prove one-hot decode against the generated map
constants, the outstanding limit and that no response is lost or
duplicated. \textbf{Neither formal job proves liveness}: nothing says a
transfer ever completes (\dcite{39}{sections 5.3, 5.5}).

On the one recipe every block here is measured with --- Yosys
\code{stat -liberty} on \code{ihp-sg13g2} typical, \code{abc} at a 20\,ns
target, one gate equivalent $=$ 7.2576\,$\mu$m$^2$ (\dcite{45}{section 6})
--- the fabric is 868
cells, 55 flip-flops, \textbf{10{,}928.8872\,$\mu$m$^2$}, 3.96\,\% of the
unprotected core, and the APB bridge 201 cells, 94 flip-flops,
6{,}437.4912\,$\mu$m$^2$ (\dcite{60}{section 5.2}).

\section{The memory map, and its single source}

\begin{sloppypar}
\code{regmap/memmap.yaml} \emph{is} the memory map. It is not a table in
prose: \code{regmap/generate\_memmap.py} emits the reader's document, the
Python constants the test suites use, the RTL's address decoder, both
device tables' contents, the C header the bare-metal programs include and
the linker's \code{MEMORY} block, and \code{sw/tests/test\_memmap.py}
fails if any is stale (\dcite{39}{section 2.1}).
\end{sloppypar}

\begin{sloppypar}
Three rules are enforced by the generator rather than commented. Ibex resets to \code{\{boot\_addr\_i[31:8], 8'h80\}}, so the generator
derives the reset vector and refuses a boot region that does not contain
it or cannot execute it. Ibex's \code{mtvec} is
256-byte aligned and vectored-only, so \textbf{every region base must be
256-byte aligned} and the map is never itself why a trap vector lands
elsewhere. And every region is a naturally aligned power of two with no
overlaps, so decode is a prefix compare rather than two magnitude
comparisons in the address path of every access.
\end{sloppypar}

\begin{table}[tbp]\small
\setlength{\tabcolsep}{4pt}
\centering
\caption{The system-bus address map, from \code{regmap/memmap.yaml} and
its generated form \code{docs/memmap-soc.md}. Ten regions, six
implemented. Everything outside every region reaches the error slave.}
\label{tab:arch-memmap}
\begin{tabular}{@{}lllllp{5.1cm}@{}}
\toprule
Base & Size & Name & Port & Status & What it is \\
\midrule
\code{0x00000000} & 32\,KiB   & RAM    & \code{ram}   & implemented & System SRAM, SECDED per byte lane, scrubbed \\
\code{0x10000000} & 256\,MiB  & NPU    & \code{npu}   & implemented & Accelerator window; sixteen 4\,KiB node windows \\
\code{0xC0000000} & 8\,KiB    & ROM    & \code{rom}   & implemented & Boot ROM; reset vector at \code{0xC0000080} \\
\code{0xD0000000} & 32\,MiB   & QSPI3  & ---          & reserved    & Flash execute-in-place, 3-byte addresses \\
\code{0xD8000000} & 128\,MiB  & QSPI4  & ---          & reserved    & Flash execute-in-place, 4-byte addresses \\
\code{0xE0000000} & 64\,KiB   & CLINT  & \code{clint} & implemented & \code{msip}, \code{mtimecmp}, \code{mtime} \\
\code{0xF8000000} & 4\,MiB    & PLIC   & ---          & reserved    & Interrupt controller; \cref{sec:arch-irq} \\
\code{0xFE000000} & 16\,MiB   & DEBUG  & ---          & reserved    & RISC-V debug module; there is none \\
\code{0xFF900000} & 1\,MiB    & APB    & \code{apb}   & implemented & Peripheral bridge, sixteen 4\,KiB slots \\
\code{0xFFFFF000} & 4\,KiB    & PNP    & \code{pnp}   & implemented & System-bus device table \\
\bottomrule
\end{tabular}
\end{table}

\begin{sloppypar}
\Cref{tab:arch-memmap} is the map, and every base is taken unchanged from
\dcite{08}{section 3.1}, which took them from the reference family ---
except the NPU window, the one row that document marks
project-specific.  The RAM is \textbf{32\,KiB,
not 64}, since the memory-protection work: the region shrank rather than
the macros growing, and \cref{sec:arch-mem} is why.
\dcite{60}{section 5.3} carries 64\,KiB in its revision-0.1 body with a
dated correction beside it; the map is the normative source.
\end{sloppypar}

\section{The memories}
\label{sec:arch-mem}

There are two implementations of one module name, and which one a reader
is looking at decides what the numbers mean. \code{soc\_mem.v} is
behavioural and is what every simulation and formal job reads;
\code{soc\_mem\_sram.v} is a drop-in with the same module name and port
list built on real vendor macros, and \textbf{no simulation ever reads
it}.

\begin{sloppypar}
The macro arrangement, confirmed from \code{soc\_mem\_sram.v} and
\dcite{67}{sections 3, 8}, is \textbf{eight macros}: four
\code{RM\_IHPSG13\_1P\_2048x64\_c2\_bm\_bist} for the RAM, the smallest
byte-writable 64\,KiB build the PDK~\cite{ihppdk} admits (the denser \code{8192x32} part has no byte-mask port, and Ibex emits
\code{sb} and \code{sh}); two \code{RM\_IHPSG13\_1P\_1024x32\_c2\_bm\_bist} for the
ROM, because there is no \code{2048x32} part; and two
\code{RM\_IHPSG13\_1P\_512x16\_c2\_bm\_bist}, one per ROM bank, holding the (39,32) word code's seven check bits at
$+90{,}618.62\,\mu$m$^2$. The check pair was designed unplaced and was \textbf{placed and routed on
2026-09-13} in run \code{s83romecc5}, with zero detailed-routing DRC
errors and a die that does not grow by one database unit. The estimate
that the die would have to grow by about 200\,$\mu$m in $x$ to take the
pair was \textbf{superseded on 2026-09-13} and is kept rather than
deleted: the space it needed was vertical, in the pocket the shorter ROM
macro leaves in its own row, so the estimate had measured the wrong axis.
\end{sloppypar}

\subsection{Why the check bits cost half the RAM}

A \code{2048x64} row is 64 bits and holds two 32-bit words. The project's
SECDED code is (72,64) --- eight bits \emph{wider} than the row --- and a
(39,32) word code needs fourteen spare bits per row. \textbf{There is no
bit-widening in a hard macro}, so the check field goes into another macro
or into the space a narrower word leaves. \dcite{67}{section 3} costs six
options from the vendor LEF, and the three that keep 64\,KiB all buy it
with a read-modify-write: a seventh \code{2048x64} for the (39,32) word
code at $+491{,}633.62\,\mu$m$^2$, $+21.88$\,\% of the six-macro area,
with one on every \code{sb} and \code{sh}; one \code{4096x16} for a
(72,64) code over the whole row at $+257{,}608.51\,\mu$m$^2$,
$+11.46$\,\%, with one on \emph{every} store; and four more
\code{2048x64} at $+1{,}966{,}534.46\,\mu$m$^2$ to avoid it. That is a design
change and not a wrapper: a second port cycle per store, a hazard against
the pipelined request the fabric permits, and a grant the memory would
have to withhold.

What was built holds \textbf{one 32-bit word per 64-bit row as four
(16,8) byte codewords}, which is exactly 64 bits. Each is the frozen
(72,64) Hsiao code~\cite{hsiao1970} with data bits 8--63 tied to zero; three of the eight
check rows are then identically zero, so what is stored is a (13,8) code
plus three constant bits, and (13,8) is the minimal SECDED width for a
byte. A byte write lands through the macro's own per-bit write mask with
its check bits beside it: no seventh macro, no die growth, no
read-modify-write, and the response stays one cycle.  What the 32\,KiB that went were holding was measured rather than assumed:
the 28-check bring-up program uses \textbf{192 bytes} of static data plus
a stack, and the RTL keeps the two-words-per-row arm for the day a sized
workload wants 64\,KiB back (\dcite{67}{section 3.4}).

Every read decodes the row on the way out: a single-bit error in any lane
is corrected before the read data and counted, and an uncorrectable
syndrome answers with a bus error that Ibex takes as an access fault, so
\textbf{the recovery is a trap and not a wrong word}. A scrubber~\cite{saleh1990} sits
behind the row port rather than on the fabric --- a third bus master would
have been a re-proof of the arbiter and its ownership queues --- reading
one row in an idle cycle, examining it in the next, and writing back only
the lanes the decoder corrected. \textbf{The bus always wins}: if the bus
asks in the examine cycle the scrub result is dropped, so a bus write
between a scrub read and its write-back cannot be overwritten.

The protection costs \textbf{$+31{,}391.3124\,\mu$m$^2$, $+4.783$\,\%}
on the whole SoC at the SRAM boundary, taking it from 43{,}884 cells and
5{,}599 flip-flops to 45{,}952 and 5{,}833 (\dcite{67}{section 5.1}).
Routed, the read-return group's worst slack at
\code{nom\_slow\_1p08V\_125C} falls from $-2.7112$ to
$-4.0021$\,ns, and it is still not the binding path. Two limits close the
section: the row address into
either memory is \textbf{unprotected}, since a corrupted index reads the
wrong row, which is a valid codeword; and on this PDK the boot ROM is an
SRAM that powers up undefined and nothing in the design can write it, so
check bits on a ROM nobody can write are check bits nobody can initialise
(\dcite{67}{sections 8, 9}).

\section{Time, interrupts and supervision}

The CLINT and the general-purpose timer are both timers and are not
duplication. \textbf{The CLINT owns time}: \code{mtime} is a free-running
64-bit up-counter that is never reloaded and never restarted, and it is
the only time base with an architectural meaning. \textbf{The GPTIMER
owns intervals}: its timers count down from a reload value and restart,
which is what periodic work and timeouts want, and its last timer is the
watchdog. Neither is a convenient way to do the other's job, and the two
have independent time bases that nothing synchronises
(\dcite{40}{section 4}).

The CLINT sits on the system bus rather than behind the APB bridge, which
costs $+1{,}421.92\,\mu$m$^2$, $+0.52$\,\% of the core --- the difference
of two measurements, and tagged an estimate where it is published
(\dcite{40}{section 9}) --- and buys a
one-cycle \code{mtime} read instead of the bridge's three-cycle minimum
on the one register software polls in a loop. It has the standard layout
--- \code{msip} at \code{0x0000}, \code{mtimecmp} at \code{0x4000},
\code{mtime} at \code{0xBFF8} --- one hart, and a bus error at every other
offset. One behaviour is deliberately not papered over: the naive two-store
64-bit \code{mtimecmp} update glitches, and the cocotb suite
\emph{asserts} that a spurious interrupt occurs on it rather than only
checking the safe sequence.

\code{mtime} is the one block here whose hardening is a code rather than
a copy. \dcite{58}{} makes its 64 stored bits the data field of a (72,64)
SECDED codeword, decoded and corrected on every tick, at \textbf{eight
added flip-flops} and $+7{,}026.8310\,\mu$m$^2$. It corrects: 144 of 144 single-bit injections into the 72 stored bits came
back corrected, against an unprotected design that
displaced the clock in 95 of 128 draws and announced none. That
document's own negative result sits beside it --- a residue-mod-3
detector measures \textbf{98.694\,\% of the corrector's area} and
corrects nothing --- and \code{mtimecmp} gets nothing, because 128
injections into it displaced the clock zero times.

The CLINT's area is published twice and both numbers are right:
\dcite{40}{section 9} reports 1{,}038 cells and 16{,}683.48\,$\mu$m$^2$,
\dcite{45}{section 6.2} reports 1{,}043 and 16{,}712.8164 for the same
unchanged file, and the difference is the source list the script read on
the day. The second is the one to quote for the SoC as it stands.

\subsection{No platform interrupt controller}
\label{sec:arch-irq}

Thirteen interrupt sources are defined and Ibex has fifteen dedicated,
individually vectored, individually maskable, priority-ordered fast local
interrupt inputs, each with its own vector-table entry. That is an
interrupt controller, in the core, at no cost. \textbf{No PLIC is
built}: the region stays reserved and faults, \code{irq\_external\_i} is
tied low and left free so that adding a controller later is a wiring
change, and \textbf{lines 13 and 14 are unassigned}
(\dcite{40}{section 3}). The condition that reverses the decision is
mechanical: the generator refuses a line outside 0--14 and refuses two
sources on one line, so the first peripheral for which no free line
remains is a build failure naming the decision, and
\code{sw/tests/test\_memmap.py} recomputes the two spare lines from the
map rather than asserting them.

The map keeps \textbf{two interrupt namespaces} that must not be
conflated: \code{irq} is the plug-and-play source number in the device
table's identification word, running to 24, and \code{line} is the Ibex
wire index, running to 12.  Four core inputs are architectural rather than
per-peripheral: MSOFT (ID 3) and MTIMER (ID 7) from the CLINT, MEXT
(ID 11) unconnected, and \textbf{NMI (ID 31) from the watchdog}.

\subsection{The three-stage watchdog}

The watchdog is not a timer with a reset on the end. It carries the
weight the \code{SecureIbex} decision put on it --- it is the only
backstop a silently corrupted processor has --- and its specification is
numbered properties in its own file header, against which the
formal properties and the cocotb suite are written
(\dcite{40}{section 5}, \dcite{60}{section 5.6}). There are nine of them:
\dcite{60}{section 5.6} tabulates eight, and \code{soc\_wdog.v} carries a
ninth, W9, added when \dcite{75}{} found the reset request reaching an
asynchronous reset through a combinational decode of the voted word.

\textbf{W1: armed at reset, not disableable by software.} The enable
reads 1 out of power-on reset and a write of 0 is accepted by the bus and
\emph{ignored}, so a driver that tried can discover that it failed. Only
a bootstrap pin sampled once at power-on reset release can stop it, and
the status register reports what was sampled. \textbf{W2: the time base is not software-programmable}, because the
attack on W1 is not to disable the watchdog but to slow it: on a stock
GPTIMER it shares a writable prescaler. Here it has its own fixed divider that no
register reaches, and the longest timeout it can be talked into is a
constant of the netlist, $2^{16}\times16 = 1{,}048{,}576$ clocks. There
is deliberately no clamp register: \textbf{a bound enforced by the width
of a counter cannot be misconfigured and one enforced by a comparator
can}, and the first version's comparator turned out to be dead logic
pretending to be a defence.

\textbf{W3: three stages.} Stage 1 raises a non-maskable interrupt and
reloads --- non-maskable because the core clears \code{mstatus.MIE} on
entry to every trap handler, so a core stuck inside a handler would never
take a maskable warning in exactly the case it exists for. Stage 2, on a second expiry with stage 1 unacknowledged, asserts the
system reset, which the core cannot prevent. Stage 3, after two watchdog
resets, latches an external pin, because resetting has demonstrably not
worked and the decision belongs outside the chip. \textbf{W4: the state
survives the reset it causes}, so the whole block is in the power-on
domain and the SoC has two reset domains with a synchroniser. \textbf{W5:
keyed writes}, upper half-word \code{0xA51F}, on every register and not
just the kick, because an unkeyed acknowledge would let a core spraying
stores hold the machine in stage 1 for ever. \textbf{W6:} the
persistent, silent state is bundled into one word, tripled and voted
every cycle~\cite{johnson2010}, while the flip-flops the block rewrites
every tick or kick are left alone. \textbf{W7 and W8}, a windowed kick mode and a kick
budget, are both measured; W7 is \textbf{recommended for removal}. The
word was 22 bits when \dcite{41}{} measured it, against 36 unprotected;
W7's and W8's four fields joined it by the same criterion, so the shipped
block protects \textbf{29 bits} against 44
(\code{hw/soc/rtl/soc\_wdog.v}, \dcite{75}{section 3.5}).

``Armed at reset'' was right and not sufficient, and running it is what
showed that. The reload register is in the power-on domain by W4, so a
short timeout software installed before it went wrong \emph{survived the
reset that going wrong caused}, and the SoC looped with the console never reaching its first character; a
stage-2 reset now restores the reload to the maximum. \textbf{Putting
state outside the reset domain is what makes the record survive and also
what makes a stale configuration survive} (\dcite{40}{section 7.2}). The
redundancy itself was counted rather than assumed, after stripping every
protective attribute~\cite{esa_ftmr} from the source text: \textbf{102 mapped
flip-flops} on the build \dcite{41}{} measured, 36 unprotected plus three
replicas of 22, at $+5{,}091.06\,\mu$m$^2$ and \textbf{1.6\,\% of what
lockstep would have cost}; \textbf{132} on the block as it stands,
$44 + 1 + 3 \times 29$, the one being W9's registered request
(\dcite{75}{section 3.5}).

\section{The peripheral bus and what is on it}

Sixteen 4\,KiB slots sit inside the 1\,MiB APB window. The bridge
\textbf{does not decode which peripheral is selected}: it drives one ``some peripheral'' select and
\code{soc\_top.v} turns that into per-slot selects from the generated
constants, because a bridge that decoded slots would be a second copy of
the memory map, and the second copy is where drift lives
(\dcite{39}{section 4}).

\begin{table}[tbp]\small
\centering
\caption{The peripheral slots, from \code{regmap/memmap.yaml}. Areas are
the same recipe as every other block here, rounded here to two decimals,
and they come from different sessions: UART0, TIMER0, BUSSTAT and APBPNP
from \dcite{45}{section 6.2}, GPIO from \dcite{65}{section 7}, QSPI from
\dcite{66}{section 9}, SCRUB from \dcite{67}{section 5.1} and BOOTREG ---
the protected build, which is what the design instantiates --- from
\dcite{69}{section 7}. A block re-measured in a different source list
moves by of order one per cent (\dcite{45}{section 4.3}), and the BUSSTAT
row is the file \dcite{44}{} measured: \dcite{55}{} later added the NPU
fault lines at a further $+4{,}796.1396\,\mu$m$^2$, which
\dcite{60}{section 5.7} declines to fold in because no run has reported
the sum.}
\label{tab:arch-slots}
\setlength{\tabcolsep}{4pt}
\begin{tabular}{@{}llp{2.5cm}p{4.3cm}rrr@{}}
\toprule
Slot & Address & Block & Registers & IRQ & Line & Area, $\mu$m$^2$ \\
\midrule
\code{0x000} & \code{0xFF900000} & UART0, console & 4: data, status, ctrl, scaler & 2 & 0 & 4{,}359.78 \\
\code{0x001} & \code{0xFF901000} & UART1 & \emph{reserved} & 3 & 1 & --- \\
\code{0x002} & \code{0xFF902000} & GPIO, 16 pins & 10, plus 9 logical aliases & 4 & 2 & 13{,}162.87 \\
\code{0x008} & \code{0xFF908000} & TIMER0 and watchdog & scaler, reload, config, latch, 3 timers, \code{WDOGSTAT} & 8 & 3 & 33{,}528.07 \\
\code{0x009} & \code{0xFF909000} & TIMER1 & \emph{reserved} & 12 & 4 & --- \\
\code{0x00D} & \code{0xFF90D000} & SpaceWire & \emph{reserved} & 16 & 5 & --- \\
\code{0x011} & \code{0xFF911000} & CAN & \emph{reserved} & 18 & 6 & --- \\
\code{0x012} & \code{0xFF912000} & SPI & \emph{reserved} & 19 & 7 & --- \\
\code{0x013} & \code{0xFF913000} & I2C & \emph{reserved} & 20 & 8 & --- \\
\code{0x014} & \code{0xFF914000} & QSPI flash & 7, at \code{0x00}--\code{0x18} & 21 & 9 & 19{,}690.62 \\
\code{0x015} & \code{0xFF915000} & BUSSTAT & 7, at \code{0x00}--\code{0x18} & 22 & 10 & 7{,}060.55 \\
\code{0x016} & \code{0xFF916000} & SCRUB & 6 counters, 6 stickies, 2 addresses & 23 & 11 & 16{,}710.96 \\
\code{0x017} & \code{0xFF917000} & BOOTREG & 4: straps, status, report, epoch & --- & --- & 12{,}495.89 \\
\code{0x018} & \code{0xFF918000} & CLKGATE & \emph{reserved} & --- & --- & --- \\
\code{0x019} & \code{0xFF919000} & NPUCFG & 12, at \code{0x000}--\code{0x02C} & 24 & 12 & in \code{soc\_npu} \\
\code{0x0FF} & \code{0xFF9FF000} & APBPNP & two words per slot & --- & --- & 905.27 \\
\bottomrule
\end{tabular}
\end{table}

\begin{sloppypar}
\Cref{tab:arch-slots} is the current map: \textbf{nine slots implemented
and seven reserved}. \dcite{60}{section 4.1} counts nine reserved, which
was true on its own date of 2026-09-05; the SCRUB and BOOTREG slots were
implemented afterwards by \dcite{67}{} and \dcite{68}{}, and revision 0.2
did not recount them. The map is the normative source
(\dcite{60}{section 0.2}).
\end{sloppypar}

The \textbf{console UART} mirrors the reference family's register map at
four offsets and is a real serialiser, not a simulation print hook: the
testbench decodes its line rather than snooping the register write, so
the baud divider and the shift register are inside what the run proves.
It is \textbf{transmit only}: a driver that waits for the data-ready bit
waits for ever, which is correct for a part with no receiver and better
than a receiver that appears to exist.

The \textbf{GPIO} port was the first peripheral promoted from reserved to
implemented, and it was chosen first for a process reason rather than a
mission one: no block had ever walked the sequence from map promotion through
top-level decode and source lists to a bring-up check and a proof, and
walking it on this project's own RTL means that when the second block
fails a step the failure is in the block and not in the path
(\dcite{65}{section 2}). It is 16 pins on one shared interrupt line, and \textbf{its data register
reads the pad and never the output register}: a block that folded one
into the other would pass a write-then-read-back test with no pin present
at all.

The \textbf{QSPI flash controller} is register-mode with two chip
selects. The alternative, wrapping OpenTitan's \code{spi\_host}, was measured
rather than dismissed: it synthesises to
\textbf{65.051\,kGE, 171\,\% of the Ibex core}, with about four fifths of
its 5{,}404 flip-flops FIFO storage at depths fixed in a generated
package --- the area a fact, the share an estimate
(\dcite{65}{section 8.4}). What was written instead is
19{,}690.6248\,$\mu$m$^2$, \textbf{4.2\,\% of the size} of the block it
declined. Its data phase \textbf{pauses at every word}, so a transfer of any length
needs one word of storage and no FIFO --- ``the storage is where software
is, not where the block is'' (\dcite{66}{section 3.1}).

Execute-in-place is \textbf{not} built and the two windows stay reserved
with the reason in the map beside them: it would make the controller a
fabric slave as well as an APB one --- a seventh port and a re-proof of
\code{soc\_bus.v} --- and a quad read of one word is about 56 system
clocks (\dcite{66}{section 2.1}) on a core where one extra cycle of
memory latency was measured at 25.23\,\% of run time
(\dcite{50}{section 5.2}). The controller is checked against a behavioural \code{W25Q128JV} written
from the part's datasheet, with every departure from its AC table counted
and asserted zero --- after which \dcite{66}{section 12} says the
sentence it most needs to: \textbf{a behavioural model is not a
flash}.

\section{The boot flow}

Until \dcite{68}{} the boot ROM held the bring-up program, and that had
stopped working: the 28-check program was 7{,}794 of the 8{,}064 bytes
the map leaves above the reset vector when \dcite{66}{} measured it, and
7{,}826 after the memory-protection work. Four questions were open, in
four documents --- how the ROM gets contents, who writes the loader, what
the bootstrap pin and the boot report are, and what happens when boot
fails --- and they are one
design, because each answer is the next one's mechanism.

\textbf{The ROM holds a loader.} Four candidates were costed against one
criterion: do the contents exist before the core leaves reset without
needing something else with the same problem? A mask-programmed array is
the only terminal answer --- its contents exist because the mask exists,
and a bit that is a via has no stored-upset mechanism, so it needs no
check field. An SRAM loaded through the vendor macros' test port
relocates the question: \textbf{the loader then has a loader}. This PDK
ships no ROM generator~\cite{ihppdk}, so the two \code{1024x32} macros are a
\textbf{stand-in}, and in silicon on this PDK the ROM's contents are
unprotected and undefined. The loader is 3{,}104 of the 8{,}064 bytes
available.

\textbf{An uninitialised protected word is an uncorrectable, not a
zero.} The RAM's check bits do not exist until something writes them, so
the loader sweeps every word before it reads any, before the stack
pointer is set --- the only order that works, because its own stack, data
and BSS are inside what it is initialising --- and it writes \emph{whole
words}, because a byte store makes one of a row's four codewords valid
and leaves three as the SRAM powered up. \textbf{The image checksum is then computed from the RAM, not from the
flash}: a second pass reads the copied words back out and sums them, so a
word that arrived intact and landed on a row whose codeword did not take
is caught, and that pass is also where an uncorrectable in the loaded
image is detected. The sweep costs 6.003 cycles per word, the copy 44.06
and the verify 17.01.

\textbf{The escalation is the watchdog's; there is no second ladder.}
Within a boot, the primary image's magic, header checksum, geometry, body
checksum or ECC each send the loader to a secondary image elsewhere in the device. Both
failing means the loader shortens the watchdog to its minimum and stops
kicking --- safe only because a stage-2 reset restores the reload to the
maximum. The attempt limit is \textbf{three}, tied
to the watchdog's escalation count of two by an argument rather than by
taste: the external pin is asserted on the second watchdog reset, which
is the reset that begins boot number 2, and the loader's last attempt is
boot number 2, so \textbf{software gives up after the hardware has told
the platform, never before}. There is no ``wait for an upload'' state, because the console is
transmit-only. The watchdog is kicked \textbf{on progress, not on
a schedule}, from inside the copy loop gated on the word index.

\code{soc\_boot.v} is the register block and its design rule is one
sentence: \textbf{the one field with authority is the one software cannot
write}. The boot counter is maintained by hardware on each release of the
system reset, saturating, and no APB write moves it; the report and epoch
words are software's and are \emph{evidence}. The block therefore needs
no write key --- a field hardware maintains cannot be forged --- and its
only outputs are the APB read path, so it can shorten no budget and reach
no timer.

Hardening it (\dcite{69}{}) produced a ranking rather than a count. Of 92
flip-flops \textbf{eighteen are protected} --- the boot counter's eight,
the four sampled straps, the watchdog-pin readback, two flags, a two-bit
sampling delay and the
system-reset sample --- bundled into one word with a five-bit mismatch
report and tripled under the
frozen voter, while \textbf{74 are deliberately not}, 64 of them the
report and epoch words, 69.6\,\% of the block's state with zero boot
consequence~\cite{mukherjee2003}. The campaign moved one bit into the protected set: the
system-reset sample is rewritten every clock, which normally means a
register sheds an upset on its own, and that is true of the register and
false of the consequence --- an upset there manufactures a release edge
and the saturating counter never takes it back. Protection costs
$+5{,}255.9010\,\mu$m$^2$ and $+51$ flip-flops; placed and routed
(\dcite{70}{}) it costs \textbf{no die area, no hold margin and no
routing quality} --- and the layout it landed in reports
\textbf{4.30\,ns less setup slack} at \code{nom\_slow\_1p08V\_125C} than
the layout placed beside it, on a design that closes neither way, which
that document attributes to the flow's timing repair rather than to the
block: \code{u\_boot} appears on no violating path in either run.

Boot end to end, through the real fabric from a modelled flash with the
RAM powered up undefined, is \textbf{195{,}068 cycles}; the program that
follows is \textbf{220{,}256}, for \textbf{415{,}324} in the run. The
same program fetched from the boot ROM is 217{,}630, so the move costs
$+1.207$\,\%, and all of that is the fabric rather than the memory: with
the program in ROM the two masters never contend, and with it in RAM both
go to one slave port and one waits (\dcite{68}{sections 8, 8.1}).

\section{Telemetry and discovery}

Two blocks exist because a SECDED register file corrects silently, and in
silicon a corrected upset is indistinguishable from no upset at
all~\cite{santos2023}. \textbf{BUSSTAT} carries four saturating counters and four
sticky bits for register-file corrections, corrected read cycles,
uncorrectable syndromes and watchdog voter mismatches. The counters are different quantities and the header insists they not be
folded: the correction counter counts an upset the program does not
overwrite exactly once, while the corrected-read counter is an upper
bound, because one corrupt word read in a loop raises it every
iteration. The
record lives in the power-on domain and survives the reset the watchdog
causes, while the interrupt enable is deliberately in the system domain,
so an armed interrupt cannot fire on a fresh boot before software has
configured it (\dcite{60}{section 5.7}). \textbf{SCRUB} is the memories' equivalent and is a separate block for a
blunt reason --- BUSSTAT is full --- carrying the same three quantities
per memory, both scrubbers' enables, the interval, and the byte offset of
the last uncorrectable row. Both scrubbers ship \emph{walking}, by the
rule that a mechanism which ships disabled and has never caught anything
is a liability.

Two \textbf{plug-and-play device tables} are read-only in the reference
family's record layout: one for the system bus at \code{0xFFFFF000},
eight words per slave record, one for the peripheral bus at the bridge
base $+\,$\code{0xFF000}, two words per slot, both generated from the
same source as the map. The layout is mirrored because a well-known
layout is what makes a table readable by someone other than its author,
and two claims are explicitly not made: it is not an AMBA AHB
plug-and-play table, because there is no AHB, and \textbf{the device IDs
are assigned by this project and are not registered} with the vendor
whose ``various contributions'' code they sit under, so a family-aware tool shows a known vendor and an unrecognised device,
which is the truthful outcome. One field cannot express this map exactly --- the bank address register
compares \code{addr[31:20]} and four regions are smaller than 1\,MiB ---
so the exact byte size of every region is carried in user-defined word 1.

\section{What the top level is, and what it is not}

\textbf{There is no pad ring and no package}, so \code{soc\_top}'s port
list is what a floorplan or a future pad ring would have to serve. Its functional pins are the clock, the power-on reset, the watchdog
bootstrap pin, four bootstrap straps, the UART transmit line, the
watchdog's stage-3 output, sixteen GPIO pins as three wires each because
there is no pad, and the flash controller's clock, two chip selects and
four lanes as three wires each. Everything else is an observation output for a testbench or a scope. \dcite{60}{section 6.2} says the
design has three functional pins; that was true of revision 0.1, before
the GPIO, flash and strap pins.

The design has \textbf{one clock and one mode}: no clock-domain crossing,
no retention, no power domains, and a reserved address
where the clock-gate block would be. \dcite{60}{section 4.3} calls the
consequence its one measured requirement gap, four orders of magnitude
wide: on the telemetry mission profile's duty cycle the part is
0.0044\,\% busy --- an estimate, on \dcite{53}{sections 6.2, 7.3} ---
and 100\,\% powered. That second half is revision 0.1's.
Since \dcite{76}{} the netlist carries \textbf{three integrated clock
gates} where it carried one, on the core's, the fabric's and the
accelerator's clocks, the two new ones driven by their own block's
enable and the CLKGATE slot still reserved:
the share of flip-flops behind a gate goes
from 39.5\,\% to \textbf{76.0\,\%} and idle power at the typical corner
from 18.745 to \textbf{8.287\,mW}, a factor of 2.26, carried through
\dcite{60}{section 11.3}'s correction of 2026-09-12 as an
orbit average of 8.289\,mW. The gap is narrower and still orders of
magnitude wide.

Three absences belong beside \cref{tab:arch-slots}, which a reader could
otherwise take at face value. There is \textbf{no DMA, no descriptor ring and no third bus master}: a
ring is a bus master, the fabric has two hardcoded master ports and a
fairness property written for two, so a third is a fabric change plus a
re-proof --- and at this scale it would save nothing, because the
transport and not the CPU is the bottleneck
(\dcite{51}{section 12}). There is \textbf{no debug module and no JTAG}. And there are
\textbf{no spacecraft link interfaces} --- SpaceWire, CAN, SPI and I2C
are frozen addresses with nothing behind them, though all four candidate cores have been fetched, pinned and measured
and an APB-to-Wishbone bridge for two of them is proved but instantiated
nowhere (\dcite{65}{sections 9.4, 9.5}).

The whole design at the SRAM boundary, synthesised as one netlist at the
20\,ns target with both memories protected, is \textbf{45{,}952 cells,
5{,}833 flip-flops, 687{,}776.3676\,$\mu$m$^2$} of standard cells plus
the eight vendor macros (\dcite{67}{section 5.1}). That row is from
before the boot block existed --- \code{soc\_boot.v} does not appear
anywhere in \dcite{67}{} --- and the netlist the placed design of record
was built from is a different configuration, hardened boot block and
\code{ROM\_HARDEN} = 0: \textbf{45{,}546 cells, 5{,}873 flip-flops,
685{,}153.2744\,$\mu$m$^2$} (\dcite{69}{section 7.4},
\dcite{60}{section 9.10.1}). \textbf{No document measures both together},
so no combined figure is given here.
Neither number is a frequency, and none is claimed.
